
\section{Structure of elementary systems}
\SGARefTarget{sga3:xx:section:2}

\paragraph{\textbf{Theorem 2.1.}}
\SGARefTarget{sga3:xx:2.1}
\emph{Let \(S\) be a scheme and \((G,T,\alpha)\) an
\(S\)-elementary system. There exist a morphism of
\(\mathscr O_S\)-modules}
\[
  \mathfrak g^\alpha\otimes_{\mathscr O_S}\mathfrak g^{-\alpha}
    \longrightarrow\mathscr O_S,
  \qquad
  (X,Y)\longmapsto\langle X,Y\rangle,
\]
\emph{and a morphism of \(S\)-groups}
\[
  \alpha^\ast\colon\mathbb G_{m,S}\longrightarrow T
\]
\emph{such that, for every \(S'\to S\) and all}
\[
  X\in\Gamma\bigl(S',
    \mathfrak g^\alpha\otimes_{\mathscr O_S}\mathscr O_{S'}\bigr),
  \qquad
  Y\in\Gamma\bigl(S',
    \mathfrak g^{-\alpha}\otimes_{\mathscr O_S}\mathscr O_{S'}\bigr),
\]
\emph{one has}
\[
  \exp(X)\mathbin{\cdot}\exp(Y)\in\Omega(S')
  \quad\Longleftrightarrow\quad
  1+\langle X,Y\rangle\in\mathbb G_m(S').
\]
\emph{Under these conditions,}
\SGARefTarget{sga3:xx:2.1:F}
\begin{equation}\tag{F}
\begin{aligned}
  \exp(X)\mathbin{\cdot}\exp(Y)
   ={}&\exp\left(\frac{Y}{1+\langle X,Y\rangle}\right)
       \alpha^\ast\bigl(1+\langle X,Y\rangle\bigr)\\
     &{}\exp\left(\frac{X}{1+\langle X,Y\rangle}\right).
\end{aligned}
\end{equation}
\emph{Moreover, the morphisms
\((X,Y)\mapsto\langle X,Y\rangle\) and \(\alpha^\ast\) are uniquely
determined. The first is an isomorphism, and therefore places
\(\mathfrak g^\alpha\) and \(\mathfrak g^{-\alpha}\) in duality;
furthermore,}
\[
  \alpha\circ\alpha^\ast=2,
\]
\emph{where \(2\) denotes squaring on \(\mathbb G_{m,S}\).}

In view of the uniqueness assertions, it is enough to prove the
theorem locally on \(S\). We may therefore assume that
\(\mathfrak g^\alpha\) and \(\mathfrak g^{-\alpha}\) are free on
\(S\). Choose
\[
  X\in\Gamma(S,\mathfrak g^\alpha)^\times,\qquad
  Y\in\Gamma(S,\mathfrak g^{-\alpha})^\times,
\]
and put
\[
  p_\alpha(x)=\exp(xX),\qquad
  p_{-\alpha}(y)=\exp(yY)
\]
for \(x,y\in\mathbb G_a(S')\), \(S'\to S\). By \SGARefLink{sga3:xx:1.5}{1.5} and \SGARefLink{sga3:xx:1.21}{1.21}, it
is enough to prove the following proposition.

\paragraph{\textbf{Proposition 2.2.}}
\SGARefTarget{sga3:xx:2.2}
\emph{Let \(S\) be a scheme, \(G\) an \(S\)-group, \(T\) a torus
of \(G\), and \(\alpha\) a character of \(T\) that is nontrivial
on every fiber. Let}
\[
  p_\alpha\colon\mathbb G_{a,S}\longrightarrow G,\qquad
  p_{-\alpha}\colon\mathbb G_{a,S}\longrightarrow G
\]
\emph{be group monomorphisms normalized by \(T\), with
multipliers \(\alpha\) and \(-\alpha\), respectively. Suppose
that:}
\begin{enumerate}
\item[\textup{(i)}] \emph{the morphism}
  \[
    \mathbb G_{a,S}\times_S T\times_S\mathbb G_{a,S}
      \longrightarrow G,
    \qquad
    (y,t,x)\longmapsto p_{-\alpha}(y)\,t\,p_\alpha(x),
  \]
  \emph{is an open immersion; denote its image by \(\Omega\);}
\item[\textup{(ii)}] \emph{for every \(s\in S\), the groups}
  \[
    (p_\alpha)_s\bigl(\mathbb G_{a,\kappa(s)}\bigr)
    \quad\text{and}\quad
    (p_{-\alpha})_s\bigl(\mathbb G_{a,\kappa(s)}\bigr)
  \]
  \emph{do not commute.}
\end{enumerate}
\emph{Then there exist unique elements}
\[
  a\in\mathbb G_a(S),\qquad
  \alpha^\ast\in
    \operatorname{Hom}_{S\text{-gr.}}(\mathbb G_{m,S},T)
\]
\emph{with the following properties: for every \(S'\to S\) and
all \(x,y\in\mathbb G_a(S')\),}
\[
  p_\alpha(x)p_{-\alpha}(y)\in\Omega(S')
  \quad\Longleftrightarrow\quad
  1+axy\in\mathbb G_m(S'),
\]
\emph{and, under this condition,}
\[
  p_\alpha(x)p_{-\alpha}(y)
   =p_{-\alpha}\left(\frac{y}{1+axy}\right)
      \alpha^\ast(1+axy)
      p_\alpha\left(\frac{x}{1+axy}\right).
\]
\emph{Moreover, \(a\) is invertible, that is,
\(a\in\mathbb G_m(S)\), and \(\alpha\circ\alpha^\ast=2\).}

\medskip
\noindent\emph{Proof.}
\textup{A)} Consider the morphism
\[
  \mathbb G_{a,S}^{\,2}\longrightarrow G,\qquad
  (x,y)\longmapsto p_\alpha(x)p_{-\alpha}(y).
\]
Let \(U\) be the inverse image of \(\Omega\). It is an open subset
of \(\mathbb G_{a,S}^{\,2}\) containing
\(0\times_S\mathbb G_{a,S}\) and
\(\mathbb G_{a,S}\times_S0\). There are therefore unique
morphisms of \(S\)-schemes
\[
  A\colon U\longrightarrow\mathbb G_{a,S},\qquad
  C\colon U\longrightarrow\mathbb G_{a,S},\qquad
  B\colon U\longrightarrow T
\]
satisfying the set-theoretic identity
\[
  p_\alpha(u)p_{-\alpha}(v)
    =p_{-\alpha}\bigl(A(u,v)\bigr)\,
      B(u,v)\,p_\alpha\bigl(C(u,v)\bigr).
\]
Immediately,
\[
\begin{gathered}
  A(0,v)=v,\qquad A(u,0)=0,\qquad
  C(u,0)=u,\qquad C(0,v)=0,\\
  B(u,0)=B(0,v)=e.
\end{gathered}
\]

Let \(S'\) be a separated \(S\)-scheme and let \(t\in T(S')\).
Since \(\Omega_{S'}\) is stable under \(\intaut(t)\), the last
formula of \SGARefLink{sga3:xx:1.16}{1.16} shows that \(U_{S'}\) is stable under the
automorphism
\[
  (x,y)\longmapsto
    \bigl(\alpha(t)x,\alpha(t)^{-1}y\bigr)
\]
of \(\mathbb G_{a,S'}^{\,2}\), and
\[
\begin{aligned}
 A\bigl(\alpha(t)u,\alpha(t)^{-1}v\bigr)
   &=\alpha(t)^{-1}A(u,v),\\
 C\bigl(\alpha(t)u,\alpha(t)^{-1}v\bigr)
   &=\alpha(t)C(u,v),\\
 B\bigl(\alpha(t)u,\alpha(t)^{-1}v\bigr)
   &=B(u,v).
\end{aligned}
\]
Because \(\alpha\) is faithfully flat, it follows that for every
\(S'\to S\) and every \(z\in\mathbb G_m(S')\), the open subset
\(U_{S'}\) is stable under
\((x,y)\mapsto(zx,z^{-1}y)\), and
\[
\begin{aligned}
  A(zu,z^{-1}v)&=z^{-1}A(u,v),\\
  C(zu,z^{-1}v)&=zC(u,v),\\
  B(zu,z^{-1}v)&=B(u,v).
\end{aligned}
\]

Suppose first that \(v\) is invertible. Taking \(z=v\), we see
that if \((u,v)\) is a section of \(U\), then \((uv,1)\) is also
a section of \(U\), and
\[
  A(uv,1)=v^{-1}A(u,v),\qquad
  B(uv,1)=B(u,v).
\]
Let \(V\) be the open subset of \(\mathbb G_{a,S}^{\,2}\) defined
by\footnote{The editors added the condition
``\((1,uv)\in U(S')\)''.}
\[
  (u,v)\in V(S')
  \quad\Longleftrightarrow\quad
  (u,v),\ (uv,1),\ \text{and }(1,uv)\text{ belong to }U(S').
\]

Since \(U\) contains
\(0\times_S\mathbb G_{a,S}\) and
\(\mathbb G_{a,S}\times_S0\), the open subset \(V\) is a
neighborhood of the zero section of \(\mathbb G_{a,S}^{\,2}\).
The preceding calculation shows that, on
\(V\cap(\mathbb G_{a,S}\times_S\mathbb G_{m,S})\), the two
morphisms
\[
  (u,v)\longmapsto A(u,v),\qquad
  (u,v)\longmapsto vA(uv,1)
\]
coincide, as do
\[
  (u,v)\longmapsto B(u,v),\qquad
  (u,v)\longmapsto B(uv,1).
\]
Since
\(\mathbb G_{a,S}\times_S\mathbb G_{m,S}\) is schematically dense
in \(\mathbb G_{a,S}^{\,2}\), these morphisms coincide on \(V\).

We have \(A(0,1)=1\). Hence there is an open subset \(W_1\) of
\(\mathbb G_{a,S}\), containing the zero section, such that
\(A(x,1)\) is invertible for every section \(x\) of \(W_1\). Put
\[
  A(x,1)^{-1}=F(x).
\]
If \((u,v)\in V(S')\)\footnote{Here and below, the editors replaced
\(U(S')\) by \(V(S')\).} and \(uv\in W_1(S')\), then
\[
  A(u,v)=vA(uv,1)=vF(uv)^{-1}.
\]
The same argument applied to \(C\) gives an open subset \(W_2\)
of \(\mathbb G_{a,S}\), containing the zero section, and an
element \(E\in\mathscr O(W_2)^\times\)\footnote{The editors use
\(E\) for the element denoted \(G\) in the original, since \(G\)
already denotes the \(S\)-group under consideration.} such that
\[
  C(u,v)=uC(1,uv)=uE(uv)^{-1}
\]
whenever \((u,v)\in V(S')\) and \(uv\in W_2(S')\).

Consequently, with \(W=W_1\cap W_2\), there are \(S\)-morphisms
\[
\begin{array}{lll}
  F\colon W\to\mathbb G_{m,S},&\qquad&F(0)=1,\\
  H\colon W\to T,&&H(0)=e,\\
  E\colon W\to\mathbb G_{m,S},&&E(0)=1,
\end{array}
\]
such that, if \((u,v)\in V(S')\) and \(uv\in W(S')\), then
\SGARefTarget{sga3:xx:2.2:plus}
\begin{equation}\tag{+}
  p_\alpha(u)p_{-\alpha}(v)
    =p_{-\alpha}\bigl(vF(uv)^{-1}\bigr)\,
       H(uv)\,p_\alpha\bigl(uE(uv)^{-1}\bigr).
\end{equation}
